\section{Coverage}
\label{section:node-level:coverage}

%Gibt es ein Amdahl law?
%   \item[$\boxplus$] Use Amdahl's law to characterise your code. If it fits to
%   the law, i.e.~there is an $f_{\text{intra}}$ such that
%   \eqref{equation:amdahl:speedup}.


Studying code is an option. 

Ratio there are two options: number of lines covered vs runtime ratio


Zoom into time where only one thread is available
Filter out other time
Filter out other threads
Look at the top code at this time
Study if it is parallelised



Linaro gives you an overview of the toatl OpenMP time in the summary
The overall performance report gives you that information as well, but it
doesn't break it down over time




Intel APS is another tool where you see to which degree we are covered by
parallelism.


Once we have zoomed in, we can look if this code should actually be covered by
parallelism




